%% ---------------------------------------------------------------------
\section{¿Por qué empezar por una dimensión?}
\label{sec:un-porque}
%% ---------------------------------------------------------------------

Este capítulo estudia el caso más sencillo posible: un sistema invariante en el tiempo, sin entrada, con una sola variable de estado,
\begin{equation}
  \dot x=f(x),\qquad x(0)=x_0,\qquad x\in\R .
  \label{eq:un-escalar}
\end{equation}
La razón para empezar aquí no es que el caso escalar sea útil por sí mismo, aunque lo es (el diodo túnel, una población con cosecha y el juego halcón--paloma de este capítulo son escalares), sino que en una dimensión \emph{todo se ve}. Las ideas que ocupan el resto del curso aparecen aquí sin el ruido de la geometría del plano: equilibrio, estabilidad local, linealización y sus límites, región de atracción, función potencial (un primer encuentro con las funciones de Lyapunov), conjunto invariante, bifurcación, existencia y unicidad, escape en tiempo finito. Cada una se volverá a encontrar, generalizada.

\notamargen{\textbf{Ruta de las clases.} Clase~1: \S\ref{sec:un-flujo}--\ref{sec:un-monotonia}. Clase~2: \S\ref{sec:un-poblaciones}--\ref{sec:un-escape}. Las dos últimas secciones motivan el \cref{cap:preliminares}.}

\begin{pausa}
Antes de leer, apueste: (1) ¿Puede oscilar un sistema $\dot x=f(x)$ con $x\in\R$? (2) Si $f(\eq)=0$ y $f'(\eq)=0$, ¿qué se concluye sobre $\eq$? (3) ¿Cuánto vale en $t=2$ la solución de $\dot x=x^2$, $x(0)=1$? (4) ¿Pueden dos soluciones distintas de $\dot x=f(x)$ salir del mismo punto?
\respuesta{(1) No, si $f$ es localmente Lipschitz: las soluciones son monótonas (\cref{prop:un-monotonia}). Sí, si el sistema es variante en el tiempo o si el estado vive en una circunferencia (\cref{tr:un-oscilar}). (2) Nada: puede ser estable, inestable o semiestable (\cref{tr:un-nohiperbolico}). (3) No existe: la solución escapó a infinito en $t=1$ (\cref{ej:un-escape}). (4) Sí, si $f$ no es Lipschitz, como en $\dot x=x^{1/3}$ (\cref{ej:un-raiz}); no, si lo es.}
\end{pausa}
